\documentclass[12pt,a4paper]{article}

% Documento compatible con pdflatex.
\usepackage[utf8]{inputenc}
\usepackage[T1]{fontenc}
\usepackage[spanish,es-nodecimaldot]{babel}
\usepackage[a4paper,margin=2.5cm]{geometry}
\usepackage{graphicx}
\usepackage{float}
\usepackage{amsmath,amssymb}
\usepackage{booktabs}
\usepackage{array}
\usepackage{enumitem}
\usepackage{xcolor}
\usepackage{hyperref}
\usepackage{fancyhdr}
\usepackage{longtable}
\usepackage{url}

\hypersetup{colorlinks=true,linkcolor=blue!50!black,urlcolor=blue!50!black,citecolor=blue!50!black,pdftitle={Sistema educativo de iluminación y balance energético hotelero}}
\setlength{\parindent}{0pt}
\setlength{\parskip}{0.35em}
\setlist{nosep,leftmargin=*}
\renewcommand{\arraystretch}{1.08}
\pagestyle{fancy}
\fancyhf{}
\lhead{Iluminación y balance energético hotelero}
\rhead{Ciencias de la Computación 1}
\cfoot{\thepage}

\newcommand{\clase}[1]{\texttt{#1}}
\newcommand{\estado}[1]{\texttt{#1}}
\newcommand{\codigo}[1]{\texttt{#1}}

\begin{document}

\begin{titlepage}
    \centering
    \vspace*{1.5cm}
    {\Large \textbf{Ciencias de la Computación 1}\par}
    \vspace{0.9cm}
    {\Huge \textbf{Sistema educativo de iluminación}\par}
    {\Huge \textbf{y balance energético para un hotel}\par}
    \vspace{1.2cm}
    {\Large Documento de arquitectura y propuesta de implementación\par}
    \vfill
    \begin{tabular}{rl}
        \textbf{Escenario:} & Hotel de dos torres en Bogotá\\
        \textbf{Actividades:} & Habitaciones, cocina, recreación y parqueadero\\
        \textbf{Análisis:} & Modelo unificado de 24 horas\\
        \textbf{Implementación propuesta:} & Python\\
    \end{tabular}
    \vfill
    {\large \today\par}
\end{titlepage}

\tableofcontents
\newpage

El documento se organiza únicamente según los contenidos de evaluación solicitados. El escenario es un hotel compuesto por exactamente dos torres en Bogotá, con cuatro actividades y un análisis temporal único de 24 horas. Las cantidades, rangos y geometría son datos sintéticos educativos; el sistema diagnostica y recomienda, pero no constituye una certificación fotométrica, eléctrica ni arquitectónica.

\section{Entradas}

Las entradas son los datos que el programa necesita para construir el escenario, calcular la iluminación y ejecutar el balance energético.

\textbf{Almacenamiento de las entradas.} La información inicial se guarda en el archivo \codigo{datos\_hotel.json}. Este archivo contiene las torres, pisos, zonas, ventanas, vidrios, paredes, luminarias, huéspedes y sus restricciones, horarios de actividad y de cierre, perfiles climáticos y parámetros de evaluación. Python carga el archivo y convierte los registros en objetos \codigo{dataclass}. Además, se crean diccionarios indexados por identificador para localizar rápidamente habitaciones, zonas, ventanas y huéspedes. Las mismas entradas se pueden editar desde los formularios de la interfaz web.

\begin{center}\small
\begin{tabular}{p{0.25\textwidth}p{0.28\textwidth}p{0.34\textwidth}}
\toprule
\textbf{Información} & \textbf{Almacenamiento inicial} & \textbf{Uso en el algoritmo}\\
\midrule
Torres, pisos y zonas & \codigo{datos\_hotel.json} & Construir el árbol \clase{Hotel}--\clase{Torre}--\clase{Piso}--\clase{ZonaHotel}.\\
Ventanas, paredes y luminarias & \codigo{datos\_hotel.json} & Calcular sombras, visibilidad, reflexión y luz artificial.\\
Vidrios y reflectancias de pared y de fachada & \codigo{datos\_hotel.json} & Aplicar transmitancia, refracción, reflexión interior y reflexión exterior entre torres.\\
Clima y horarios & \codigo{datos\_hotel.json} & Calcular la luz base y el modo de cada zona en cada instante del eje de 24 horas.\\
Huéspedes, asignaciones y restricciones & \codigo{datos\_hotel.json} & Determinar ocupación, prioridades y posibles reasignaciones.\\
Parámetros de evaluación & \codigo{datos\_hotel.json} & Nivel de emergencia, umbrales de habitabilidad, pesos de reasignación y rangos de reestructuración.\\
Resultados temporales & Memoria de Python, índice rojo-negro y \codigo{resultados.json} & Clasificar zonas, consultar franjas horarias y generar el reporte final.\\
\bottomrule
\end{tabular}
\end{center}

El archivo de entrada conserva los datos del escenario; el Octree, el grafo y los árboles AVL y rojo-negro no se escriben manualmente. Después de cargar el JSON, el programa construye el árbol jerárquico, el Octree y el grafo de influencia. Los resultados se mantienen durante la simulación, se indexan en el árbol rojo-negro y se exportan a \codigo{resultados.json} y a un reporte CSV.

\textbf{Relación entre entradas y evaluación.} La geometría, las posiciones y las normales permiten calcular sombras y visibilidad; las ventanas, el clima y la orientación permiten calcular luz natural; el vidrio permite calcular transmisión y refracción; las paredes permiten calcular la reflexión interior y las fachadas y el grafo, la reflexión entre torres; las luminarias y sus horarios permiten calcular luz artificial; y la actividad, el modo y el perfil permiten comparar el resultado con los niveles mínimo, objetivo y máximo. Finalmente, la ocupación y la capacidad del generador determinan las prioridades del balance y la posibilidad de reasignar huéspedes.

\begin{center}\small
\fbox{\parbox{0.88\textwidth}{\centering
\codigo{datos\_hotel.json}
$\longrightarrow$ objetos \codigo{dataclass} y diccionarios por ID
$\longrightarrow$ árbol jerárquico + Octree + grafo
$\longrightarrow$ simulación de iluminación
$\longrightarrow$ \codigo{resultados.json} y reporte CSV}}
\end{center}

\textbf{Hotel y geometría.} El hotel contiene exactamente dos torres. Cada torre contiene pisos; cada piso contiene zonas; una zona puede contener ventanas, superficies y luminarias. Las áreas de referencia se inspiran en modelos comerciales del DOE \cite{doeReferenceBuildings}, pero las cajas y registros siguientes son sintéticos.

\begin{center}\small
\begin{tabular}{p{0.17\textwidth}p{0.21\textwidth}p{0.30\textwidth}p{0.23\textwidth}}
\toprule
\textbf{Torre} & \textbf{Referencia} & \textbf{Caja espacial (m)} & \textbf{Zonas y ocupación}\\
\midrule
\codigo{TP} (Torre Pequeña) & 4.013 m$^2$, 4 pisos & $x=0..32$, $y=0..31.35$, $z=0..12$ & 32 habitaciones, 4 cocinas, 4 recreaciones, 1 parqueadero; 24 ocupadas.\\
\codigo{TG} (Torre Grande) & 11.346 m$^2$, 6 pisos & $x=45..89$, $y=0..43$, $z=0..18$ & 72 habitaciones, 6 cocinas, 6 recreaciones, 2 parqueaderos; 54 ocupadas.\\
\bottomrule
\end{tabular}
\end{center}

Cada piso tiene 3 m de altura y la separación horizontal aproximada entre torres es de 13 m. Las habitaciones tienen 24 m$^2$ y capacidad para dos huéspedes. El escenario inicia con 78 huéspedes: 24 habitaciones ocupadas en la torre pequeña y 54 en la grande. Las torres se identifican como \codigo{TP} (Torre Pequeña) y \codigo{TG} (Torre Grande); el nombre completo se guarda como campo de presentación. Las zonas usan el prefijo de su torre, la actividad y un número: \codigo{TP-H-101} es una habitación de la torre pequeña y \codigo{TG-COC-01} una cocina de la torre grande (\codigo{H} habitación, \codigo{COC} cocina, \codigo{REC} recreación, \codigo{PAR} parqueadero). Los huéspedes son \codigo{H-001}--\codigo{H-078}; no representan personas reales. La geometría de las torres es fija: el sistema no mueve ni rediseña los edificios.

\textbf{Zonas y actividades.} El sistema usa cuatro actividades para evitar clases innecesarias:

\begin{center}\small
\begin{tabular}{p{0.16\textwidth}r r r p{0.19\textwidth}p{0.24\textwidth}}
\toprule
\textbf{Actividad} & \textbf{Área} & \textbf{Zonas TP} & \textbf{Zonas TG} & \textbf{Ventanas} & \textbf{Horario sintético}\\
\midrule
Habitación & 24 m$^2$ & 32 & 72 & 1 de 2,0 m$^2$ & 06:00--22:00 activa; resto reducida. Nunca cierra.\\
Cocina & 35 m$^2$ & 4 & 6 & 2 de 1,5 m$^2$ & 06:00--22:00 activa; 23:00--05:00 cerrada; resto reducida.\\
Recreación & 80 m$^2$ & 4 & 6 & 3 de 2,0 m$^2$ & 08:00--23:00 activa; 00:00--07:00 cerrada; resto reducida.\\
Parqueadero & 450 m$^2$ & 1 & 2 & Sin ventana interior & Activo las 24 horas. Nunca cierra.\\
\bottomrule
\end{tabular}
\end{center}

Los horarios se guardan en \codigo{horarios} del JSON con los campos \codigo{inicio\_activa}, \codigo{fin\_activa}, \codigo{cierre\_inicio} y \codigo{cierre\_fin}. Si una hora cae a la vez en la franja activa y en la de cierre, manda la franja activa; así la reestructuración puede extender el horario de la cocina hasta las 24:00 sin editar el cierre.

Las ventanas de habitaciones alternan fachadas norte, este, sur y oeste; las de cocinas y recreaciones miran hacia fachadas exteriores. Cada luminaria equivalente usa 100 lm/W, $CU=0,70$ y $MF=0,80$ como supuestos de cálculo.

\textbf{Perfiles de iluminación.} Los siguientes valores son presets educativos configurables, expresados como lux de referencia. No son una transcripción normativa del RETILAP:

\begin{center}\small
\begin{tabular}{p{0.20\textwidth}r r r r p{0.22\textwidth}}
\toprule
\textbf{Actividad} & \textbf{Mín. (lux)} & \textbf{Objetivo (lux)} & \textbf{Máx. (lux)} & \textbf{Prioridad} & \textbf{Luminaria}\\
\midrule
Habitación activa & 100 & 200 & 300 & 3 & 72 W\\
Cocina activa & 300 & 400 & 500 & 3 & 200 W\\
Recreación activa & 200 & 300 & 500 & 2 & 450 W\\
Parqueadero activo & 50 & 100 & 150 & 1 & 720 W\\
\bottomrule
\end{tabular}
\end{center}

En modo \estado{REDUCIDA} los tres niveles se multiplican por \codigo{factor\_reducido}, que vale 0,30 en todas las actividades. Una habitación en descanso se modela como \estado{REDUCIDA}; no se marca automáticamente como deficitaria. El nivel de emergencia es un parámetro separado, \codigo{parametros\_evaluacion.emergencia.nivel\_lux}, con valor de 10 lux (si falta en el JSON, se usa 10). Es un supuesto del proyecto cercano al promedio inicial de 1 pie-candela (10,8 lux) que NFPA 101 pide en las rutas de evacuación \cite{nfpa101}.

\textbf{Clima de Bogotá.} La ubicación de referencia es latitud $4.7110$, longitud $-74.0721$, con climatología NASA POWER de 2001--2020 \cite{nasaPowerApi}. Se usan estos perfiles:

\begin{center}\small
\begin{tabular}{p{0.22\textwidth}p{0.16\textwidth}r r r r}
\toprule
\textbf{Perfil} & \textbf{Referencia} & \textbf{Radiación} & \textbf{Nubosidad} & \textbf{Lluvia} & \textbf{factor\_luz}\\
 & & kWh/m$^2$/día & \% & mm/día & \\
\midrule
Nubes parciales & Febrero & 4,863 & 70,89 & 2,51 & 0,689\\
Nublado & Mayo & 4,477 & 85,82 & 6,68 & 0,634\\
Lluvioso & Abril & 4,454 & 84,87 & 6,97 & 0,631\\
Cielo despejado & Clear-sky & 7,063 & -- & -- & 1,000\\
\bottomrule
\end{tabular}
\end{center}

Como contexto anual se usan 4,639 kWh/m$^2$/día de radiación global all-sky, 7,063 de cielo despejado, 78,41\% de nubosidad y 4,45 mm/día de precipitación. El escenario base usa el perfil de nubes parciales. La irradiancia se transforma en lux de referencia mediante una constante calibrable, no en lux fotométricos medidos:
\begin{equation}
 E_{base}(t)=K_{lux}\,C(t)\,factor_{luz},\qquad
 factor_{luz}=\min\left(1,\frac{I_{clima}}{I_{despejado}}\right),\qquad
 C(t)=\sin\left(\frac{\pi(t-6)}{12}\right),
 \label{eq:clima}
\end{equation}
donde $K_{lux}=900$ (\codigo{parametros\_evaluacion.luz.lux\_por\_irradiancia}), $I_{clima}$ es la radiación del perfil activo e $I_{despejado}=7,063$ kWh/m$^2$/día es la radiación de cielo despejado de NASA POWER (\codigo{parametros\_evaluacion.luz.radiacion\_cielo\_despejado}; si falta, se usa la mayor radiación de los perfiles). La curva solar $C(t)$ solo se evalúa entre 06:00 y 18:00; fuera de ese intervalo $E_{base}=0$. El factor no se escribe a mano en el JSON: lo calcula la propiedad \codigo{PerfilClimatico.factor\_luz}, así que cambiar la radiación cambia el resultado. La nubosidad $c$ reparte la luz entre directa y difusa:
\begin{equation}
 f_{directa}=\max(0.10,1-c),\qquad f_{difusa}=0.15+0.45c.
 \label{eq:nubosidad}
\end{equation}
La precipitación se carga y se reporta, pero es informativa: no entra en el cálculo.

\textbf{Vidrios y paredes.} Las ventanas reciben uno de tres presets basados en fichas de vidrio arquitectónico \cite{cardinalSolarControl,cardinalClear}: \clase{VIDRIO\_CLARO}, VLT 0,82; \clase{VIDRIO\_LOW\_E}, VLT 0,70--0,72; o \clase{VIDRIO\_CONTROL\_SOLAR}, VLT 0,52. Los tres usan índice de refracción $n=1,5$ (\codigo{indice\_refraccion}), que entra en el cálculo de Snell y Fresnel de la luz directa. El vidrio, el color de pared y la orientación de las ventanas se configuran por actividad en \codigo{acabados} y por recinto en \codigo{zonas\_personalizadas}. Las paredes usan reflectancias educativas: blanco 0,80, crema 0,70, gris claro 0,60, gris medio 0,45 y oscuro 0,25. Cada torre tiene además cuatro fachadas (norte, sur, este y oeste) cuya reflectancia sale del campo \codigo{fachada} de la torre, con los mismos presets; por defecto es \codigo{GRIS\_CLARO} (0,60).

\textbf{Fuentes y generador.} Una fuente contiene posición, potencia, eficacia, $CU$, $MF$, horario y zona servida. El generador tiene capacidad nominal sintética de 22,0 kW y se usa solamente para iluminación. No alimenta climatización, ascensores ni otras cargas. La reserva del 20\% es el margen con que se dimensionó la capacidad (ecuación~\ref{eq:generador}); en cada instante el balance puede repartir los 22,0 kW completos.

\textbf{Huéspedes y restricciones.} Cada huésped puede tener restricciones en \codigo{huespedes.restricciones\_por\_huesped}. Hay dos restricciones duras, que descartan candidatas: \codigo{MISMA\_TORRE} y \codigo{PISO\_MAXIMO:n}. Hay una preferencia, que suma al puntaje de la candidata: \codigo{ORIENTACION:ESTE}, \codigo{NORTE}, \codigo{OESTE} o \codigo{SUR}. Una restricción desconocida se rechaza al cargar el escenario.

\section{Salidas}

El programa produce resultados por zona e instante y después los agrega por piso, torre, actividad y hotel. La salida temporal conserva separadas las componentes para que la luz artificial no se confunda con iluminación natural.

\begin{center}\small
\begin{tabular}{p{0.25\textwidth}p{0.63\textwidth}}
\toprule
\textbf{Salida} & \textbf{Contenido}\\
\midrule
\clase{ResultadoIluminacion} & Hora, periodo, modo, luz natural, luz reflejada, luz artificial, luz total, potencia asignada y estado.\\
Estado por zona & Déficit, confort o exceso; actividad, perfil y causa dominante.\\
Veredicto de habitabilidad & \estado{HABITABLE}, \estado{CASI\_HABITABLE} o \estado{NO\_HABITABLE} por zona, con puntaje total, diurno y nocturno, brecha, causa y recomendación; puntaje ponderado del hotel.\\
Balance & Potencia máxima, capacidad del generador, demanda mínima, demanda objetivo, energía en Wh y kWh.\\
Reporte por franja horaria & Para el día completo y las franjas 00--06, 06--12, 12--18, 18--22 y 22--24: energía en Wh y kWh, energía por actividad y por torre, hora y potencia pico del generador, déficit residual (lux$\cdot$h bajo el mínimo), energía no asignada (Wh que habrían cerrado ese déficit), zonas críticas (las 10 con más horas en déficit), horas con déficit por zona, instantes con déficit y estado predominante de cada zona.\\
Reporte jerárquico & Energía por zona, piso, torre y hotel; comparación entre torres.\\
Reasignación & Por huésped en habitación \estado{NO\_HABITABLE}: habitación propuesta y su puntaje, candidatas revisadas y descartadas con motivo, o motivo de falta de alternativa.\\
Mejoras y reestructuración & Cambio medido del puntaje para cada mejora candidata y secuencia de cambios de la reestructuración.\\
Estructuras de datos & Pares evaluados y aristas del grafo, tamaño, altura, rotaciones y comparaciones del AVL y del rojo-negro.\\
\bottomrule
\end{tabular}
\end{center}

Las salidas se exportan a \codigo{resultados.json}, \codigo{resumen.json} y \codigo{reporte.csv} con \codigo{python -m src.main}, y se muestran en la interfaz web.

Para cualquier instante se reporta:
\begin{equation}
 E_{total}(z,t)=E_{natural}(z,t)+E_{reflejada}(z,t)+E_{artificial}(z,t).
 \label{eq:totalhotel}
\end{equation}
La salida distingue además puntuación diurna/nocturna, potencia y energía. No se suman lux con Wh: son magnitudes distintas.

\section{Condiciones}

Las condiciones definen cuándo se calcula cada componente y cómo se interpreta el resultado.

\textbf{Condición temporal.} El eje cubre de 00:00 a 24:00 con un único conjunto de 48 instantes separados por $\Delta t=0,5$ h. Con 127 zonas, cada simulación produce 6.096 resultados. El periodo es \codigo{DIA} entre 07:00 y 17:00, \codigo{TRANSICION} entre 06:00 y 07:00 y entre 17:00 y 18:00, y \codigo{NOCHE} el resto. Durante el día y la transición se consideran luz natural, reflexión y artificial; durante la noche:
\begin{equation}
 E_{natural}(z,t)=0,\qquad E_{reflejada\ solar}(z,t)=0,\qquad E_{total}(z,t)=E_{artificial}(z,t).
 \label{eq:noche}
\end{equation}
No se crean clases separadas de día y noche.

\textbf{Modos de operación.} Toda zona puede estar \estado{ACTIVA}, \estado{REDUCIDA}, \estado{CERRADA} o \estado{EMERGENCIA}. Los tres primeros salen del horario; el cuarto lo asigna el balance:

\begin{center}\small
\begin{tabular}{p{0.20\textwidth}p{0.66\textwidth}}
\toprule
\textbf{Modo} & \textbf{Cuándo y efecto}\\
\midrule
\estado{ACTIVA} & Dentro de la franja activa de la actividad. Usa el perfil normal. El parqueadero está siempre activo.\\
\estado{REDUCIDA} & Fuera de la franja activa y fuera del cierre. Mínimo, objetivo y máximo se multiplican por 0,30.\\
\estado{CERRADA} & Cocina de 23:00 a 05:00 y recreación de 00:00 a 07:00; habitaciones y parqueadero nunca cierran. La zona no recibe energía y sus pasos quedan como \estado{DATOS\_INSUFICIENTES}, así que no cuentan en el puntaje de habitabilidad.\\
\estado{EMERGENCIA} & La zona abierta no alcanzó su mínimo porque el generador se agotó. Conserva la potencia que alcanzó a recibir, empezando por la reserva del nivel de emergencia de 10 lux, no compite por objetivos, se clasifica contra el mínimo del modo que le tocaba por horario y ese paso cuenta como déficit.\\
\bottomrule
\end{tabular}
\end{center}

En el escenario base el generador alcanza y ninguna zona entra en \estado{EMERGENCIA}; hay 260 pasos \estado{CERRADA}. Con el generador de 10 kW del caso \codigo{deficit\_leve} aparecen 16 pasos \estado{EMERGENCIA}, y con el de 4 kW del caso \codigo{deficit\_severo}, 667. ISO 30061:2007 y NFPA 101 se usan como referencia de contexto para emergencia \cite{isoEmergency,nfpa101}.

\textbf{Clasificación por instante.} El estado siempre usa el nivel real de iluminación de la zona, no su peso de prioridad:

\begin{center}\small
\begin{tabular}{p{0.29\textwidth}p{0.55\textwidth}}
\toprule
\textbf{Estado} & \textbf{Regla}\\
\midrule
\estado{DEFICIT\_SEVERO} & $E<0,50E_{min}$.\\
\estado{DEFICIT\_LEVE} & $0,50E_{min}\leq E<E_{min}$.\\
\estado{CONFORTABLE} & $E_{min}\leq E\leq E_{max}$.\\
\estado{EXCESO\_LEVE} & $E_{max}<E\leq1,25E_{max}$.\\
\estado{EXCESO\_SEVERO} & $E>1,25E_{max}$.\\
\estado{DATOS\_INSUFICIENTES} & Falta geometría, actividad, perfil o asignación, o la zona está \estado{CERRADA}.\\
\bottomrule
\end{tabular}
\end{center}

$E_{min}$ y $E_{max}$ son los del modo de la zona en ese instante: en \estado{REDUCIDA} ya vienen multiplicados por 0,30. Los factores 0,50 y 1,25 son reglas del proyecto, no valores RETILAP. Los rangos deben contrastarse con la categoría exacta, plano de trabajo, uniformidad, mantenimiento y deslumbramiento antes de un diseño real \cite{retilap,en12464}.

\textbf{Veredicto de habitabilidad por zona.} Después de las 24 horas, cada zona recibe un veredicto a partir de sus $n$ pasos evaluables (los pasos \estado{DATOS\_INSUFICIENTES}, incluidos los cerrados, no cuentan):
\begin{equation}
 P(z)=\frac{1}{n}\sum_{t}q(t),\qquad
 q(t)=\begin{cases}
 1 & \text{si el estado es \estado{CONFORTABLE}},\\
 0,5 & \text{si es \estado{DEFICIT\_LEVE} o \estado{EXCESO\_LEVE}},\\
 0 & \text{si es \estado{DEFICIT\_SEVERO} o \estado{EXCESO\_SEVERO}}.
 \end{cases}
 \label{eq:puntaje}
\end{equation}

\begin{center}\small
\begin{tabular}{p{0.29\textwidth}p{0.55\textwidth}}
\toprule
\textbf{Veredicto} & \textbf{Regla}\\
\midrule
\estado{NO\_HABITABLE} & $P<0,75$ o más del 5\% de los pasos en \estado{DEFICIT\_SEVERO}.\\
\estado{HABITABLE} & $P\geq0,90$ y no es \estado{NO\_HABITABLE}.\\
\estado{CASI\_HABITABLE} & $0,75\leq P<0,90$ y no es \estado{NO\_HABITABLE}.\\
\estado{DATOS\_INSUFICIENTES} & La zona no tiene pasos evaluables.\\
\bottomrule
\end{tabular}
\end{center}

El puntaje se calcula también solo con los pasos de día y solo con los de noche. Para una zona casi habitable se reporta la brecha hasta 0,90, la causa más frecuente (déficit diurno o nocturno, exceso solar o artificial) y una recomendación. El puntaje del hotel es el promedio de los puntajes de zona ponderado por $Prioridad_{actividad}\,(1+0,10(p-1))$. Los umbrales y los valores de $q$ están en \codigo{parametros\_evaluacion} del JSON. En el escenario base 101 zonas son \estado{HABITABLE}, 26 quedan \estado{CASI\_HABITABLE}, todas por exceso solar, ninguna es \estado{NO\_HABITABLE} y el puntaje ponderado del hotel es 0,9726 (día 0,9451, noche 1,0).

\textbf{Condiciones ópticas.} Para la luz directa, cada ventana calcula el ángulo de refracción con la ley de Snell y la transmitancia de Fresnel no polarizada \cite{edmundOptics}. Con $\cos\theta_i=\max(0,\mathbf n_v\cdot\mathbf s)$, donde $\mathbf s$ es la dirección solar:
\begin{equation}
 n_{aire}\sin(\theta_i)=n_{vidrio}\sin(\theta_t),\qquad
 R_s=\left(\frac{\cos\theta_i-n\cos\theta_t}{\cos\theta_i+n\cos\theta_t}\right)^2,\qquad
 R_p=\left(\frac{\cos\theta_t-n\cos\theta_i}{\cos\theta_t+n\cos\theta_i}\right)^2,
 \label{eq:snell}
\end{equation}
\begin{equation}
 T(\theta_i)=1-\frac{R_s+R_p}{2},\qquad
 F(\theta_i)=\left(\frac{T(\theta_i)}{T(0)}\right)^2,
 \label{eq:fresnel}
\end{equation}
donde $n=n_{vidrio}=1,5$ y $n_{aire}=1$. El cuadrado representa las dos caras del vidrio. Al dividir por la incidencia perpendicular, $F(0)=1$ y la VLT del catálogo se conserva a 0°. El factor angular resultante es 0,997 a 30°, 0,979 a 45°, 0,900 a 60°, 0,605 a 75° y 0,163 a 85°. La función que lo calcula es \codigo{factor\_angular\_vidrio} en \codigo{geometria.py}. La luz natural de la zona es el promedio de sus ventanas:
\begin{equation}
 E_{natural}(z,t)=\frac{1}{|V_z|}\sum_{v\in V_z}E_{base}(t)\,\tau_v\left(\cos\theta_i\,F(\theta_i)\,f_{directa}+f_{difusa}\right),
 \label{eq:natural}
\end{equation}
donde $\tau_v$ es la VLT del catálogo y $\cos\theta_i=0$ si un rayo del Octree hacia el sol encuentra la otra torre. La componente difusa usa solo la VLT, sin corrección angular. No se modelan polarización de salida, dispersión espectral, reflexión múltiple dentro del vidrio ni marcos detallados.

La otra torre también tapa parte del cielo. Por eso $f_{difusa}$ se multiplica, en cada ventana, por $1-O_v$:
\begin{equation}
 O_v=\min\left(1,\sum_{f}\cos\theta_{v,f}\,\frac{A_f}{\pi d_{vf}^{2}+A_f}\right),
 \label{eq:obstruccion}
\end{equation}
donde la suma recorre las fachadas $f$ de la otra torre que la ventana tiene enfrente, $A_f$ es el área de la fachada, $d_{vf}$ la distancia y $\cos\theta_{v,f}$ el coseno entre la normal de la ventana y la dirección hacia la fachada. En el escenario base 40 ventanas tienen obstrucción; \codigo{TP-H-101-V1}, al este y frente a TG, tiene $O_v=0,593$.

\textbf{Luz reflejada.} Sólo existe de día y suma dos partes: la reflexión exterior entre torres, que se calcula con el grafo de influencia (ecuación~\ref{eq:influencia}), y la reflexión interior del recinto. La reflexión interior usa el método de flujo dividido:
\begin{equation}
 E_{ri}(z,t)=\frac{\Phi(z,t)\,\rho_z}{A_{int}(z)\,(1-\rho_z)},\qquad
 \Phi(z,t)=\sum_{v\in V_z}\left(E_v(t)+E_{ext}(v,t)\right)A_v,\qquad
 A_{int}(z)=2A_z+4\sqrt{A_z}\cdot 3\ \text{m},
 \label{eq:interior}
\end{equation}
donde $E_v$ es la luz natural de la ventana, $E_{ext}$ la luz exterior reflejada que entra por ella, $A_v$ el área de la ventana y $\rho_z$ la reflectancia de la pared del recinto. Una habitación de 24 m² tiene $A_{int}\approx106,8$ m²; con pared blanca ($\rho_z=0,80$) el factor $\rho_z/(1-\rho_z)$ vale 4. La fórmula promedia todos los rebotes interiores, así que no es de primer orden, y hace que el color de la pared cambie la luz de su propio cuarto. La reflejada de la zona es
\begin{equation}
 E_{reflejada}(z,t)=\frac{1}{|V_z|}\sum_{v\in V_z}E_{ext}(v,t)+E_{ri}(z,t).
 \label{eq:reflejada}
\end{equation}
A las 12:00 del escenario base, \codigo{TP-H-101} (ventana al este, frente a TG) recibe 97,0 lux naturales y 98,4 reflejados; \codigo{TP-H-102} (al norte, sin torre enfrente) recibe 358,7 naturales y 26,9 reflejados, casi todos interiores; \codigo{TG-H-101} recibe 238,3 y 17,9.

La luz artificial no recibe reflexión aparte: el $CU=0,70$ del método de lúmenes ya incluye el rebote interior de la luz de la luminaria. El $CU$ es fijo y no cambia con el color de la pared.

\section{Diagrama de clases inicial y estructuras de datos}

La arquitectura inicial se representa mediante composición:
\begin{center}
\fbox{\parbox{0.82\textwidth}{\centering
\clase{Hotel} $\rightarrow$ \clase{Torre} $\rightarrow$ \clase{Piso} $\rightarrow$ \clase{ZonaHotel} $\rightarrow$ \clase{Ventana}}}
\end{center}

El diagrama de la Figura~\ref{fig:uml-clases} presenta las clases agrupadas en cinco bloques: dominio físico, energía y huéspedes, perfiles, estructuras espaciales y control. Las relaciones se representan con composición para el árbol físico y los resultados de zona, herencia para \clase{HabitacionHotel} como especialización de \clase{ZonaHotel}, asociación para el vidrio, el perfil, el generador, las fuentes, el clima y la asignación, y dependencias punteadas para los servicios de \clase{Octree}, \clase{GrafoInfluencia}, \clase{BalanceEnergetico}, \clase{EvaluadorHabitabilidad} y \clase{ReasignadorHuespedes}.

\textbf{Clases de dominio y control.}

\begin{center}\small
\begin{tabular}{p{0.29\textwidth}p{0.57\textwidth}}
\toprule
\textbf{Clase} & \textbf{Responsabilidad y relación}\\
\midrule
\clase{Hotel} & Contiene exactamente dos \clase{Torre}, un \clase{GeneradorEnergia}, perfiles, fuentes y asignaciones.\\
\clase{Torre}, \clase{Piso} & Mantienen caja espacial, pisos, peso $1+0,10(p-1)$ y zonas.\\
\clase{Fachada} & Cara exterior de una torre (cuatro por torre): lado, normal, reflectancia, caja y área. Es el origen de las aristas del grafo.\\
\clase{ZonaHotel}, \clase{HabitacionHotel} & Mantienen actividad, área, modo, perfil, ventanas, capacidad, ocupación y resultados.\\
\clase{Ventana}, \clase{Vidrio}, \clase{SuperficiePared} & Representan posición, orientación, VLT, refracción, color y reflectancia.\\
\clase{FuenteIluminacion}, \clase{GeneradorEnergia} & Mantienen potencia, eficacia, horarios, capacidad y asignación.\\
\clase{Huesped}, \clase{Asignacion} & Representan restricciones y relación huésped--habitación.\\
\clase{PerfilClimatico}, \clase{PerfilIluminacion} & Mantienen clima, mínimos, objetivos, máximos y prioridades.\\
\clase{ResultadoIluminacion}, \clase{EvaluadorHabitabilidad} & Guardan componentes temporales, clasifican zonas y emiten el veredicto de habitabilidad.\\
\clase{BalanceEnergetico}, \clase{ReasignadorHuespedes} & Distribuyen potencia y proponen habitaciones alternativas.\\
\clase{SimuladorUnificado} & Recorre los 48 instantes, construye el Octree y el grafo, y coordina balance, evaluación y reasignación.\\
\clase{Octree}, \clase{GrafoInfluencia} & Consultas espaciales de rayo y radio; aristas fachada--ventana de la otra torre indexadas por origen y por destino, y obstrucción del cielo por ventana.\\
\clase{IndiceAVL}, \clase{IndiceRojoNegro}, \clase{IndiceTemporal} & Bolsa de habitaciones libres para la reasignación e índice de resultados por hora.\\
\bottomrule
\end{tabular}
\end{center}

En el código, \clase{BalanceEnergetico}, \clase{EvaluadorHabitabilidad} y \clase{ReasignadorHuespedes} se implementan como funciones de módulo: \codigo{balancear} en \codigo{balance.py}, \codigo{evaluar\_hotel} en \codigo{habitabilidad.py} y \codigo{proponer\_reasignaciones} en \codigo{reasignacion.py}.

\textbf{Estructura para modelar edificios.} Se usará un árbol de composición mediante objetos Python o `dataclasses`: cada \clase{Hotel} contiene una lista de torres; cada torre una lista de pisos; cada piso una lista de zonas; y cada zona listas de ventanas, superficies y fuentes. Además, diccionarios indexados por ID permiten localizar rápidamente \codigo{TP-H-101}, zonas y huéspedes.

\textbf{Octree para el espacio.} Las torres, fachadas, paredes, ventanas y luminarias se representan mediante cajas alineadas con los ejes o puntos. Un \clase{Octree} almacena objetos en nodos con límites, objetos y hasta ocho hijos. Un objeto que no cabe completamente en un hijo permanece en su padre. Cada nodo admite hasta 16 objetos antes de dividirse, con profundidad máxima 6, y la raíz cubre todos los objetos del escenario con un margen de 2 m. Se usan dos consultas: \codigo{consultar\_rayo}, para saber si la otra torre hace sombra sobre una ventana o tapa la línea entre una fachada y una ventana, y \codigo{consultar\_radio}, para encontrar las fachadas de la otra torre cercanas a cada ventana al construir el grafo.

\textbf{Grafo para la reflexión entre torres.} La reflexión exterior se representa mediante un grafo dirigido y ponderado con listas de adyacencia:
\begin{center}\textit{\clase{Fachada} $\longrightarrow$ \clase{Ventana} de la otra torre}\end{center}
Los orígenes son las 8 fachadas (norte, sur, este y oeste de cada torre). Su reflectancia sale del campo \codigo{fachada} de cada torre en el JSON; si falta, se usa \codigo{GRIS\_CLARO} (0,60). Para cada ventana $v$ se toma su punto de salida en la fachada de su propia torre (\codigo{punto\_en\_fachada}); desde ahí, \codigo{consultar\_radio} (18 m) entrega las fachadas candidatas $f$ de la otra torre, y $d_{vf}$ es la distancia al punto más cercano de cada una. Se exige que ventana y fachada se miren y que un rayo entre ellas no cruce otra torre. El peso y la luz exterior que entra por la ventana son:
\begin{equation}
 w_{fv}=\rho_f\cos\theta_f\cos\theta_v\frac{A_f}{\pi d_{vf}^{2}+A_f},\qquad
 E_{ext}(v,t)=\tau_v\sum_{f\rightarrow v}w_{fv}E_f(t),
 \label{eq:influencia}
\end{equation}
\begin{equation}
 E_f(t)=E_{base}(t)\left(\max(0,\mathbf n_f\cdot\mathbf s)f_{directa}+f_{difusa}\right),
 \label{eq:fachada}
\end{equation}
donde $\cos\theta_f$ y $\cos\theta_v$ miden cuánto se miran la fachada y la ventana, $A_f/(\pi d_{vf}^{2}+A_f)$ es un factor de vista aproximado para una superficie grande y solo se guardan aristas con $w_{fv}>0,0001$. La reflexión exterior es de primer orden: la luz reflejada por una fachada no vuelve a propagarse. A las 12:00 del escenario base, \codigo{TG-F-OESTE} recibe 290,6 lux y su arista hacia \codigo{TP-H-101-V1} pesa 0,3559. Al construir el grafo también se acumula la obstrucción del cielo de cada ventana (ecuación~\ref{eq:obstruccion}).

El grafo se construye con \codigo{Octree.consultar\_radio} desde cada ventana, en vez de comparar cada fachada con todas las ventanas. Las aristas se guardan en dos diccionarios, por origen y por destino; así, en cada instante, el cálculo de la luz exterior de una ventana solo recorre las aristas que llegan a ella. En el escenario base hay 8 fachadas y 154 ventanas: la versión ingenua evaluaría 1.232 pares y la versión con Octree evalúa 68 candidatos, visitando 7.834 nodos del Octree. El grafo tiene 46 aristas: \codigo{TP-F-ESTE} llega a 24 ventanas de TG, \codigo{TG-F-OESTE} a 16 ventanas de TP y \codigo{TP-F-NORTE} a 6 ventanas de recreaciones de TG. Los resultados de iluminación son idénticos a los de la construcción ingenua; solo cambia el número de comparaciones.

\textbf{Árbol AVL para la reasignación.} La reasignación usa un \clase{IndiceAVL} como bolsa de habitaciones con veredicto \estado{HABITABLE} y al menos una cama libre. La clave es $(torre, piso, -calidad, id)$, con
\begin{equation}
 calidad(r)=w_l\,P(r)+w_e\,Ahorro(r),\qquad Ahorro(r)=1-\min\left(1,\frac{Energia(r)}{Energia_{max}}\right),
 \label{eq:calidad}
\end{equation}
donde $P(r)$ es el puntaje de habitabilidad de la habitación y $Energia_{max}$ la mayor energía consumida por una zona. Como la calidad entra negativa, dentro de cada piso la primera clave es la mejor habitación. El árbol usa cuatro operaciones: insertar, \codigo{sucesor} (menor clave $\geq x$), \codigo{siguiente} (menor clave $>x$) y \codigo{eliminar} con rebalanceo. Para cada huésped que hay que mover se consulta el sucesor de $(torre, piso, -\infty)$ en cada piso permitido por sus restricciones; si tiene preferencia de orientación, se siguen los sucesores del piso mientras alguno todavía pueda ganar con el bono. Cuando una habitación se llena, se elimina del árbol, de modo que dos huéspedes no reciben la misma cama. El trabajo es casi todo de lectura (varias consultas y como máximo una eliminación por huésped), y por eso conviene el equilibrio estricto del AVL. El costo es $O(G\cdot F\cdot\log H)$ contra $O(G\cdot H)$ de recorrer la lista, con $G$ huéspedes por mover, $F$ pisos y $H$ habitaciones.

En el escenario base el AVL se construye con las 78 habitaciones habitables con cama libre (altura 7, 58 rotaciones, 417 comparaciones de construcción) y no se consulta, porque ningún huésped está en una habitación \estado{NO\_HABITABLE}. En una prueba de carga sobre el escenario base con 60 de las 78 habitaciones ocupadas marcadas \estado{NO\_HABITABLE}, el AVL hizo 2.087 comparaciones contra 6.240 de la búsqueda lineal, con 34 eliminaciones y ninguna habitación con sobrecupo; con restricciones mezcladas, 2.045 contra 6.240. En las dos corridas la habitación elegida coincidió, huésped por huésped, con una búsqueda lineal que descuenta capacidad.

\textbf{Árbol rojo-negro para los resultados temporales.} Cada \clase{ResultadoIluminacion} se inserta, apenas se calcula, en un \clase{IndiceRojoNegro} con clave $(hora, zona\_id)$. Las franjas horarias se consultan como rango semiabierto $[inicio, fin)$ en $O(\log R+k)$, con $R$ resultados y $k$ resultados en la franja, contra $O(R)$ de recorrer la lista; construir el índice cuesta $O(R\log R)$. De esas consultas salen los reportes por franja de la sección de salidas y el selector de franja horaria del plano en la interfaz. Se eligió el rojo-negro porque hay 6.096 inserciones y pocas consultas, y en teoría rota menos que el AVL al insertar. La medición, hecha insertando los mismos resultados en los dos árboles, da:

\begin{center}\small
\begin{tabular}{p{0.52\textwidth}r}
\toprule
\textbf{Medida} & \textbf{Escenario base}\\
 & 6.096 resultados\\
\midrule
Rotaciones rojo-negro / AVL & 9.099 / 7.232\\
Altura rojo-negro / AVL & 22 / 14\\
Nodos visitados en la franja 18--22 / lista completa & 1.029 / 6.096\\
\bottomrule
\end{tabular}
\end{center}

Los casos de prueba usan el hotel completo y producen las mismas 6.096 claves, así que dan los mismos valores.

En el escenario base el rojo-negro rota más que el AVL. Los resultados llegan casi ordenados (hora ascendente y, dentro de cada hora, tramos largos de identificadores ascendentes), y con esa entrada la ventaja teórica no aparece. Con 6.096 claves en orden aleatorio sí se cumple: 3.511 rotaciones del rojo-negro contra 4.176 del AVL. Se verificó que el árbol cumple las propiedades de rojo-negro, así que la diferencia no es un error de implementación. El índice cumple su papel en las consultas por franja, pero con estos datos el argumento de las rotaciones no se sostiene.

% El diagrama fuente está en diagrama-clases-hotel.mmd (Mermaid, orientación vertical TB).
% Exportarlo a PNG o PDF con ese mismo nombre. La imagen se rota 90 grados y ocupa la hoja completa.
\begin{figure}[p]
\centering
\makebox[\textwidth]{%
\includegraphics[angle=90,height=\paperwidth,width=0.96\paperheight,keepaspectratio]{diagrama-clases-hotel.png}}
\caption{Diagrama de clases inicial del sistema hotelero.}
\label{fig:uml-clases}
\end{figure}

\clearpage

\section{Programa y balanceos}

\textbf{Programa.} La implementación está en Python por su manejo de clases, `dataclasses`, listas, diccionarios, matemáticas vectoriales y facilidad para exportar reportes. Usa solo la biblioteca estándar. El código está organizado así:

\begin{center}\small
\begin{tabular}{p{0.25\textwidth}p{0.62\textwidth}}
\toprule
\textbf{Módulo} & \textbf{Responsabilidad}\\
\midrule
\codigo{modelo.py} & Hotel, Torre, Piso, ZonaHotel, Ventana, Vidrio, perfiles, huésped, asignación, generador y enumeraciones de actividad, modo y estado.\\
\codigo{io\_datos.py} & Carga y validación del JSON, valores por defecto y ubicación de zonas, ventanas y paredes.\\
\codigo{geometria.py} & Vector, caja, rayo, Snell y Fresnel (\codigo{factor\_angular\_vidrio}).\\
\codigo{octree.py} & Octree con consultas de rayo y de radio.\\
\codigo{influencia.py} & Grafo fachada--ventana de la otra torre construido con el Octree, visibilidad, pesos, obstrucción del cielo y reflexión exterior de primer orden.\\
\codigo{clima.py} & Curva solar, dirección solar, periodo y luz base.\\
\codigo{simulacion.py} & \clase{SimuladorUnificado}: los 48 instantes, luz natural, luz de las fachadas, reflejada (exterior e interior) y artificial, y estado de cada resultado.\\
\codigo{balance.py} & Modos por horario, demandas, orden del balance con \codigo{heapq} y agregación de energía por torre y piso.\\
\codigo{evaluacion.py} & Estados instantáneos y resumen por zona.\\
\codigo{habitabilidad.py} & Puntaje y veredicto de habitabilidad por zona y del hotel.\\
\codigo{arboles.py} & \clase{IndiceAVL} e \clase{IndiceRojoNegro} con contadores de comparaciones y rotaciones.\\
\codigo{reasignacion.py} & Restricciones de huéspedes y reasignación con la bolsa AVL.\\
\codigo{temporal.py} & Índice rojo-negro de resultados y reportes por franja horaria.\\
\codigo{sugerencias.py} & Recomendaciones por zona, generador y reasignación.\\
\codigo{mejoras.py}, \codigo{reestructuracion.py} & Medición de mejoras candidatas y búsqueda de reestructuraciones.\\
\codigo{main.py} & Ejecución por consola y exportación a JSON y CSV.\\
\codigo{api.py} & Servidor HTTP de la interfaz web (\codigo{python -m src.api}).\\
\bottomrule
\end{tabular}
\end{center}

\textbf{Iluminación artificial y balance.} Para una zona de área $A_z$:
\begin{equation}
 E_{artificial}(z,t)=\frac{\eta P_{asignada}(z,t)CU MF}{A_z},\qquad
 P_{req}(E^*)=\max\left(0,\frac{E^*-E_{natural}-E_{reflejada}}{\eta CU MF/A_z}\right),
 \label{eq:artificial}
\end{equation}
donde $E^*$ es el mínimo, el objetivo o el nivel de emergencia, $\eta$ la eficacia (100 lm/W); la potencia asignada nunca supera la potencia máxima de la luminaria de la zona. La potencia del generador y la energía acumulada son:
\begin{equation}
 P_{generador}=1,20\max_t\left(\sum_zP_{obj}(z,t)\right),\qquad
 Energia_{iluminacion}=\sum_t\sum_zP_{asignada}(z,t)\Delta t.
 \label{eq:generador}
\end{equation}
Con la demanda objetivo máxima medida al dimensionar el escenario, 18.110,7 W, la primera expresión da 21,73 kW, redondeados a los 22,0 kW del JSON. Las zonas cerradas no entran en la demanda.

En cada instante el balance sigue este orden:
\begin{enumerate}
    \item Asignar a cada zona su modo por horario (\estado{ACTIVA}, \estado{REDUCIDA} o \estado{CERRADA}) y calcular luz natural y reflejada.
    \item Calcular para cada zona abierta la potencia de emergencia (hasta $\min(10\ \text{lux}, E_{min})$), la mínima y la objetivo con la ecuación~\ref{eq:artificial}.
    \item Ordenar las zonas abiertas con una cola de prioridad (\codigo{heapq}) por categoría: primero seguridad (parqueadero), luego habitaciones ocupadas y luego el resto; dentro de cada categoría, por prioridad descendente.
    \item Recorrer ese orden y reservar el nivel de emergencia para todas las zonas abiertas.
    \item Recorrer el mismo orden y completar el mínimo de cada zona. La zona que no lo alcanza porque el generador se agotó queda en \estado{EMERGENCIA}.
    \item Repartir la potencia restante hasta el objetivo con otra cola de prioridad, ordenada solo por $Prioridad(z)$; las zonas en \estado{EMERGENCIA} no participan.
    \item Guardar estado, potencia asignada y energía; insertar el resultado en el índice rojo-negro.
\end{enumerate}
La prioridad de una zona es
\begin{equation}
 Prioridad(z)=pesoActividad(z)\,(1+0,10(p-1))\,factorOcupacion(z),
 \label{eq:prioridad}
\end{equation}
con $pesoActividad$ igual a 3 para habitación y cocina, 2 para recreación y 1 para parqueadero, $p$ el número del piso y $factorOcupacion$ igual a 1,4 en una habitación ocupada y 1,0 en las demás zonas. La cola de prioridad es suficiente porque el balance solo inserta y extrae el máximo; un árbol de búsqueda no aporta nada en esta parte.

El peso decide prioridad, pero no altera el estado físico: una zona con iluminación real menor que su mínimo continúa en déficit. La estrategia es greedy y no garantiza un óptimo global. En el escenario base el balance consume 194.432,4 Wh en las 24 horas y el pico del generador es 15.933,7 W.

\textbf{Reasignación de huéspedes.} Se reasigna a un huésped solo si su habitación tiene veredicto \estado{NO\_HABITABLE}; un paso aislado bajo el mínimo no basta. Las candidatas salen de la bolsa AVL descrita en la sección anterior, filtradas por \codigo{MISMA\_TORRE} y \codigo{PISO\_MAXIMO:n}. El puntaje de una candidata $r$ es:
\begin{equation}
 Score(r)=calidad(r)+w_a\,Orientacion(r)-w_c\frac{|p_r-p_{actual}|}{p_{max}},
 \label{eq:score}
\end{equation}
donde $calidad(r)$ es la de la ecuación~\ref{eq:calidad}, $Orientacion(r)$ vale 1 si la habitación coincide con la preferencia \codigo{ORIENTACION} del huésped o si el huésped no tiene preferencia, y 0 si no coincide, y $p_{max}$ es el piso más alto del hotel. Los pesos son $w_l=0,5$, $w_a=0,2$, $w_e=0,2$ y $w_c=0,1$ (\codigo{parametros\_evaluacion.reasignacion}). Se elige el mayor puntaje y se descuenta una cama de la habitación elegida. Cada propuesta lista las candidatas revisadas y descartadas con el motivo \codigo{MENOR\_PUNTAJE}. Si no hay candidata, el motivo es uno de estos:

\begin{center}\small
\begin{tabular}{p{0.42\textwidth}p{0.44\textwidth}}
\toprule
\textbf{Motivo} & \textbf{Cuándo}\\
\midrule
\codigo{SIN\_HABITACIONES\_HABITABLES\_LIBRES} & Al empezar no había ninguna habitación habitable con cama libre.\\
\codigo{SIN\_CAPACIDAD} & Había candidatas en los pisos permitidos, pero otros huéspedes ya ocuparon sus camas.\\
\codigo{RESTRICCIONES} & Las restricciones del huésped excluyen todos los pisos con candidatas.\\
\bottomrule
\end{tabular}
\end{center}

La operación modifica la asignación lógica, no la geometría ni una reserva real. El caso \codigo{deficit\_severo} demuestra la falta de alternativa: 71 zonas quedan \estado{NO\_HABITABLE}, solo los 3 parqueaderos quedan \estado{HABITABLE} y 25 huéspedes reciben \codigo{SIN\_HABITACIONES\_HABITABLES\_LIBRES}.

\textbf{Flujo del programa.}
\begin{verbatim}
cargar escenario y validar exactamente dos torres
construir árbol, diccionarios, Octree y grafo fachada -> ventana de la
    otra torre (consultar_radio, 18 m), con la obstrucción del cielo
para cada instante t de 00:00 a 24:00, cada 0,5 h:
    determinar modo por horario y periodo de cada zona
    calcular luz natural (sombra, obstrucción, Snell y Fresnel)
    calcular reflejada: exterior (fachadas) + interior (flujo dividido)
    reservar nivel de emergencia, luego mínimos por categoría
    distribuir la potencia restante hasta los objetivos
    clasificar, guardar ResultadoIluminacion e insertarlo en el rojo-negro
calcular reportes por franja con consultas de rango del rojo-negro
evaluar el veredicto de habitabilidad de cada zona
construir la bolsa AVL de habitaciones habitables con cama libre
para cada huésped en una habitación NO_HABITABLE:
    consultar sucesores por piso permitido, elegir y descontar capacidad
exportar resultados por zona, torre, actividad y hotel
\end{verbatim}

\textbf{Mejoras medidas y reestructuración.} La geometría de las torres es fija, pero el sistema sí cambia variables de luz y mide el efecto de cada cambio con una simulación completa. Las mejoras candidatas son usar \clase{VIDRIO\_CLARO} en las habitaciones, pintar todas las paredes de blanco, aumentar 20\% la capacidad del generador y bajar el objetivo de habitación a 180 lux; para cada una se reporta el puntaje del hotel antes y después. La reestructuración busca, dentro de los rangos de \codigo{parametros\_evaluacion.reestructuracion}, cambios por actividad en el color de pared, el vidrio, la orientación de las ventanas, el objetivo de iluminación y los horarios de cocina y recreación. Acepta un cambio solo si mejora el resultado y se detiene a las 30 evaluaciones o a las 3 rondas.

\textbf{Interfaz web.} El servidor se inicia con \codigo{python -m src.api} y sirve la carpeta \codigo{web/}. La interfaz tiene formularios para editar las entradas, un plano por torre y piso con la leyenda \codigo{TP} = Torre Pequeña y \codigo{TG} = Torre Grande, un selector de franja horaria que colorea cada zona con su estado en esa franja, los veredictos de habitabilidad, las sugerencias, las mejoras medidas, la reestructuración y una tabla de estructuras de datos y comparaciones (grafo, AVL de reasignación y rotaciones del rojo-negro contra el AVL).

\textbf{Complejidad y pruebas.} Las cotas son: sombras sin Octree $O(VBT)$; sombras con Octree $O(VT(\log B+k))$; grafo ingenuo $O(S\cdot V)$, con $S$ fachadas, y grafo con consultas de radio $O(V(\log B+k))$; agregación $O(N)$; balance con cola de prioridad $O(Z\log Z)$ por instante; reasignación $O(G\cdot F\cdot\log H)$ contra $O(G\cdot H)$; índice temporal $O(R\log R)$ para construirlo y $O(\log R+k)$ por franja. Las mediciones están en la sección de estructuras de datos. Con dos torres, el beneficio del Octree en las sombras es principalmente didáctico y organizativo; en el grafo reduce los pares evaluados de 1.232 a 68. Las pruebas automáticas están en \codigo{tests/test\_casos\_plan.py} y se ejecutan con \codigo{python -m unittest} desde la raíz del repositorio. Cubren diez casos: ventana al sol sin obstáculo, ventana sombreada por la otra torre, transmisión con cada vidrio, reflexión de primer orden sin cascada (aristas solo de una fachada a ventanas de la otra torre, peso no mayor que la reflectancia, reflejada igual a exterior más interior, menos reflejada con pared oscura que con pared blanca y reflejada nula de noche), noche sin luz natural, generador que alcanza para mínimos pero no para objetivos, habitación con déficit y alternativa, habitación con déficit sin alternativa, zona cerrada y habitación en descanso sin falso déficit, y cielo nublado contra despejado.

\textbf{Balance de alcance.} La geometría de las torres es fija: el sistema no mueve ni rediseña edificios. Sí mide mejoras y propone reestructuraciones de las variables de luz: actividades (orientación de ventanas y objetivos), acabados (pared y vidrio), horarios y generador. El proyecto no incluye óptica espectral, deslumbramiento profesional, ganancia térmica, consumo total del hotel, mallas complejas ni reservas reales. Los resultados son aproximaciones educativas, y los rangos, horarios, pesos, conteos, potencias y capacidad del generador son supuestos didácticos. La precipitación es informativa. Las fachadas se modelan como superficies opacas de un solo color: la luz que el vidrio de una ventana refleja hacia afuera se descuenta como pérdida de Fresnel y no ilumina la otra torre, aunque en una fachada con mucho vidrio ese reflejo puede ser fuerte. La reflexión exterior entre torres es de primer orden; el $CU$ de la luz artificial es fijo y no cambia con el color de la pared; las superficies interiores de cada recinto se representan con un solo valor de reflectancia; y las ventanas quedan dentro de la caja de su torre, así que para la reflexión exterior se usa su punto de salida en la fachada. RETILAP, UNE-EN 12464-1, ISO 30061 y NFPA 101 deben verificarse para una instalación real. Los estudios de energía hotelera y habitaciones sirven como contexto, no como sustituto de un cálculo eléctrico \cite{hotelReview,lbnlHotelLighting,energyStarEui}.

\section*{Referencias}
\begin{thebibliography}{99}

\bibitem{nasaPowerApi}
NASA POWER, \textit{Climatology API}. Datos de irradiancia, nubosidad y precipitación para Bogotá, periodo 2001--2020. Disponible en: \url{https://power.larc.nasa.gov/api/temporal/climatology/point}. Documentación: \url{https://power.larc.nasa.gov/docs/methodology/}.

\bibitem{retilap}
Ministerio de Minas y Energía de Colombia, \textit{Reglamento Técnico de Iluminación y Alumbrado Público -- RETILAP vigente}. La página oficial informa la modificación mediante la Resolución 40286 del 23 de junio de 2026 y la organización en cuatro libros. Disponible en: \url{https://www.minenergia.gov.co/es/misional/energia-electrica-2/reglamentos-tecnicos/reglamento-t%C3%A9cnico-de-iluminaci%C3%B3n-y-alumbrado-p%C3%BAblico-retilap/}.

\bibitem{doeReferenceBuildings}
U.S. Department of Energy, \textit{Commercial Reference Buildings}. Modelos de referencia de hotel pequeño y grande. Disponible en: \url{https://www.energy.gov/cmei/buildings/commercial-reference-buildings}.

\bibitem{energyStarEui}
ENERGY STAR, \textit{What is Energy Use Intensity (EUI)?} Disponible en: \url{https://www.energystar.gov/buildings/benchmark/understand-metrics/what-eui}.

\bibitem{hotelReview}
Rodríguez-Uribe, M. C., \textit{Energy Efficiency Indicators for Hotel Buildings}, Sustainability, 2021. Disponible en: \url{https://www.mdpi.com/2071-1050/13/4/1754}.

\bibitem{lbnlHotelLighting}
Lawrence Berkeley National Laboratory, \textit{Lighting Energy Savings Opportunities in Hotel Guestrooms}. Disponible en: \url{https://eta-publications.lbl.gov/sites/default/files/lbnl-44448.pdf}.

\bibitem{isoEmergency}
International Organization for Standardization, \textit{ISO 30061:2007 Emergency lighting}. Disponible en: \url{https://www.iso.org/standard/45801.html}.

\bibitem{nfpa101}
National Fire Protection Association, \textit{NFPA 101: Life Safety Code}. Sección 7.9.2.1, iluminación de emergencia en rutas de evacuación. Disponible en: \url{https://www.nfpa.org/codes-and-standards/nfpa-101-standard-development/101}.

\bibitem{en12464}
Asociación Española de Normalización, \textit{Iluminación de los lugares de trabajo en interiores -- UNE-EN 12464-1}. Disponible en: \url{https://revista.une.org/46/iluminacion-de-los-lugares-de-trabajo-en-interiores.html}.

\bibitem{cardinalSolarControl}
Cardinal Glass Industries, \textit{LoE2-272/270 Solar Control Glass}. Disponible en: \url{https://www.cardinalcorp.com/products/coated-glass/loe2-272-270/}.

\bibitem{cardinalClear}
Cardinal Glass Industries, \textit{Quad LoE-452+ Solar Control Glass}. Disponible en: \url{https://www.cardinalcorp.com/products/coated-glass/quad-loe-452-plus/}.

\bibitem{edmundOptics}
Edmund Optics, \textit{Introduction to Basic Ray Optics}. Disponible en: \url{https://www.edmundoptics.com/knowledge-center/application-notes/optics/introduction-to-basic-ray-optics/}.

\end{thebibliography}
\end{document}
